\section{Requirements Elicitation and Analysis}\label{requirements}

This chapter states \emph{what} KompozeR must do, independently of how it is
implemented. Requirements are derived from the stakeholders (customers,
recurring registered customers, catalog administrators) and from a small set of
representative scenarios: configuring a shelving unit, asking the chatbot for
contextual help, saving and resuming a configuration, being notified of a
price/availability change, administering the catalog, and consulting order
trends. They are grouped into \emph{functional}, \emph{non-functional}, and
\emph{implementation} requirements, each with an identifier and an acceptance
criterion used during validation (\cref{validation}).

\subsection{Glossary}\label{glossary}

\begin{description}
  \item[Component] a single catalog item (e.g.~a shelf, an upright), identified
  by a SKU and belonging to a product family.
  \item[Product family] one of the four mutually incompatible catalogs
  (Tondo, Quadro, Intelligente and Kube); components of different families cannot be combined.
  \item[Configuration] a work-in-progress or saved shelving design owned by the
  CAD context, with a state, a column plan, and a derived Bill of Materials.
  \item[BOM] Bill of Materials: the list of components (with quantities) that a
  configuration resolves to, and the basis for its price.
  \item[Subscription] a user's interest in a specific component (due to have it in the configuration), used to notify the
  user when that component changes price or availability.
  \item[Role] the authorisation level of a principal: \texttt{GUEST},
  \texttt{BASE} (authenticated customer), or \texttt{ADMIN}.
\end{description}

\subsection{Functional requirements}\label{functional-requirements}

\begin{description}
  \item[FR1 --- Visual configurator] The system shall provide a visual
  configurator to compose a shelving unit within a selected product family.
  \emph{Acceptance:} starting a session and adding compatible components yields
  a coherent configuration that can be persisted and retrieved.

  \item[FR2 --- Live preview] The system shall update the configuration preview
  after every change made in the configurator. \emph{Acceptance:} adding,
  removing, or replacing a component is reflected in the returned configuration
  snapshot.

  \item[FR3 --- Live price] The system shall update the estimated price of the
  configuration after every change. \emph{Acceptance:} the price returned for a
  configuration equals the sum derived from its current BOM and catalog prices.

  \item[FR4 --- Add to cart] The system shall allow a finalised configuration to
  be added to the cart. \emph{Acceptance:} finalising a configuration produces a
  cart whose items match the configuration BOM (SKU, quantity, unit price).

  \item[FR5 --- Contextual chatbot] The system shall provide an integrated
  chatbot answering contextual questions about component prices.
  \emph{Acceptance:} a price question about a known component returns the
  current catalog price for that component within the user's own session.

  \item[FR6 --- Save configuration] The system shall let authenticated users
  save a configuration under their profile. \emph{Acceptance:} a saved
  configuration is associated with the owner and appears in that user's list.

  \item[FR7 --- Resume configuration] The system shall let authenticated users
  list and resume previously saved configurations. \emph{Acceptance:} a user
  retrieves only their own configurations and can reopen one for editing.

  \item[FR8 --- Change notifications] The system shall notify authenticated
  users of price or availability changes affecting their active or saved
  configurations. \emph{Acceptance:} an administrative change to a subscribed
  SKU results in exactly one notification delivered to each affected user, and
  to no other user.

  \item[FR9 --- Family compatibility] The system shall prevent composing a
  configuration that mixes components of different product families.
  \emph{Acceptance:} attempting to add an incompatible component is rejected and
  the configuration is left unchanged.

  \item[FR10 --- Catalog administration] The system shall let administrators
  update catalog data, including component price and availability.
  \emph{Acceptance:} an authorised admin update is persisted and becomes visible
  to catalog reads and triggers the corresponding domain event.

  \item[FR11 --- Orders and reporting] The system shall provide administrators
  with an order view and a synthetic sales-trend report. \emph{Acceptance:} an
  authorised admin obtains order data and an aggregated trend over a selected
  period.
\end{description}

\subsection{Non-functional requirements}\label{non-functional-requirements}


\begin{description}
  \item[NFR1 --- Consistency of shared catalog] Price and availability changes
  shall be propagated consistently to every context that depends on them (cart,
  notifications). \emph{Acceptance:} after a catalog change, dependent contexts
  eventually reflect the new value, and each change is applied at most once
  (idempotent event handling).

  \item[NFR2 --- Real-time delivery] Notifications and chatbot replies shall be
  delivered to the interested user in near real time, without client polling.
  \emph{Acceptance:} a subscribed, connected user receives a push within a short
  bound of the triggering event.

  \item[NFR3 --- Fault tolerance of the core workflow] The configuration
  context shall tolerate the loss of a single database node without data loss
  and without permanent unavailability. \emph{Acceptance:} killing one replica
  member triggers automatic election of a new primary and writes resume.

  \item[NFR4 --- Recoverability] The collaborative configuration session state
  shall survive a service crash. \emph{Acceptance:} after crashing and
  restarting the service, the last checkpointed session state is recovered.

  \item[NFR5 --- Security / authorisation] Administrative functionality shall be
  accessible only to authenticated users with an authorised role, and users
  shall only access their own resources. \emph{Acceptance:} unauthenticated or
  \texttt{BASE} access to admin endpoints is rejected; cross-user access to
  configurations, orders, chat, and notifications is denied.

  \item[NFR6 --- Data minimisation / privacy] The system shall store only the
  data required to operate accounts, saved configurations, and application
  operations. \emph{Acceptance:} persisted schemas contain no data beyond what
  the use cases require.

  \item[NFR7 --- Maintainability / evolvability] Configuration, catalog,
  notifications, and reporting shall remain independently deployable and
  separately owned. \emph{Acceptance:} each context is a separate service with
  its own database and can be built, tested, and deployed independently.

  \item[NFR8 --- Usability and responsiveness] The interface shall give
  immediate feedback on relevant actions and be usable on desktop and mobile.
  \emph{Acceptance:} configuration, saving, and price updates produce visible
  feedback; core flows work at desktop and mobile breakpoints.
\end{description}

\subsection{Implementation requirements}\label{implementation-requirements}


\begin{description}
  \item[IR1 --- Distributed architecture] The system shall be realised as
  multiple, independently deployable services rather than a monolith, to
  exercise distributed-systems concerns (course constraint).

  \item[IR2 --- Containerisation] Every component shall be packaged as a
  container and be reproducibly deployable via a declarative orchestration
  (Docker Compose for development, Kubernetes for a production-like setup).

  \item[IR3 --- Technology baseline] Backend services shall be implemented in
  Node.js/TypeScript with Express, and the frontend as a Vue~3 SPA.

  \item[IR4 --- Persistence] Persistent state shall be stored in MongoDB, one
  logical database per service, honouring the database-per-service constraint.

  \item[IR5 --- Automated validation] The system shall be covered by automated
  unit, integration, and end-to-end tests runnable from the command line.
\end{description}

\subsection{Relevant Distributed System Features}
\label{ds-features}

This section motivates which distributed-system features are central to
KompozeR and which are secondary.

\paragraph{Highly relevant.}
\begin{itemize}
  \item \textbf{Resource sharing} --- \emph{central}. A single authoritative
  catalog (prices, availability) is shared by many concurrent users and by
  several contexts (configurator, cart, notifications, chatbot). Much
  of the design revolves around sharing this resource without letting contexts
  read each other's private state.

  \item \textbf{Fault tolerance and dependability} --- \emph{central} for the
  core configuration workflow. Losing a configuration is unacceptable, so its
  persistence is replicated and must survive a node failure; the collaborative
  session state must be recoverable after a crash (NFR3, NFR4). For non-core
  contexts, availability is desirable but not critical.

  \item \textbf{Consistency / data management} --- \emph{central}. Catalog
  changes must reach dependent contexts reliably and be applied exactly once
  (NFR1). The system favours \emph{eventual consistency} across contexts via
  events, while keeping strong consistency \emph{within} the core context via a
  configurable write/read concern on the replica set.

  \item \textbf{Security and trust} --- \emph{central}. The system stores
  personal and order data, distinguishes roles, and must isolate users from one
  another (NFR5). Authentication and authorisation are enforced at the gateway
  and re-checked by services.

  \item \textbf{Evolvability / maintainability} --- \emph{central}, and a
  primary reason for choosing microservices (NFR7, IR1): each bounded context
  evolves and deploys independently.
\end{itemize}

\paragraph{Relevant but bounded.}
\begin{itemize}
  \item \textbf{Scalability} --- \emph{relevant}. The workload is many
  independent, read-heavy sessions; the stateless service tier behind the
  gateway can scale horizontally. 

  \item \textbf{Transparency} --- \emph{relevant}. Location and access
  transparency are provided to clients: they only know the gateway and are
  unaware of service placement, replication of the CAD database, or which node
  is primary. 
  \item \textbf{Openness / interoperability / heterogeneity} --- \emph{partially
  relevant}. Services interoperate through plain HTTP/JSON and Socket.IO, which
  keeps them loosely coupled and language-agnostic; however the stack is
  deliberately homogeneous (Node.js), so heterogeneity is not exploited.

  \item \textbf{Performance / concurrency / efficiency} --- \emph{relevant}.
  Perceived responsiveness of preview/price updates and real-time delivery
  matter (NFR2, NFR8), and many sessions run concurrently.
\end{itemize}

\paragraph{Secondary.}
\begin{itemize}
  \item \textbf{Geographic distribution} --- the deployment is logically
  centralised (single cluster); components are distributed across
  containers/pods, not across datacenters or continents.

  \item \textbf{Economy / cost} --- beyond running on a modest, single-cluster
  footprint, cost optimisation is not a design driver in this academic setting.
\end{itemize}
